\documentclass[a4paper,11pt,twoside]{report}
% THIS FILE SHOULD BE COMPILED BY pdfLaTeX

% ----------------------   PREAMBLE PART ------------------------------

% ------------------------ ENCODING & LANGUAGES ----------------------

\usepackage[utf8]{inputenc}
%\usepackage[MeX]{polski} % Not needed unless You have a name with polish symbols or sth
\usepackage[T1]{fontenc}
\usepackage[english, polish]{babel}


\usepackage{amsmath, amsfonts, amsthm, latexsym} % MOSTLY MATHEMATICAL SYMBOLS

\usepackage[final]{pdfpages} % INPUTING TITLE PDF PAGE - GENERATE IT FIRST!
%\usepackage[backend=bibtex, style=verbose-trad2]{biblatex}


\usepackage{commath} % various commands which can make writing math expressions easier --- documentation available at: https://ctan.gust.org.pl/tex-archive/macros/latex/contrib/commath/commath.pdf

\usepackage[hidelinks]{hyperref} % for hyperlinks, for example, urls, references to equations, entries in a bibliography --- hidelinks option removes rectangles around hiperlinks


% ---------------- MARGINS, INDENTATION, LINESPREAD ------------------

\usepackage[inner=20mm, outer=20mm, bindingoffset=10mm, top=25mm, bottom=25mm]{geometry} % MARGINS


\linespread{1.5}
\allowdisplaybreaks         % ALLOWS BREAKING PAGE IN MATH MODE

\usepackage{indentfirst}    % IT MAKES THE FIRST PARAGRAPH INDENTED; NOT NEEDED
\setlength{\parindent}{5mm} % WIDTH OF AN INDENTATION


%---------------- RUNNING HEAD - CHAPTER NAMES, PAGE NUMBERS ETC. -------------------

\usepackage{fancyhdr}
\pagestyle{fancy}
\fancyhf{}
% PAGINATION: LEFT ALIGNMENT ON EVEN PAGES, RIGHT ALIGNMENT ON ODD PAGES 
\fancyfoot[LE,RO]{\thepage} 
% RIGHT HEADER: zawartość \rightmark do lewego, wewnętrznego (marginesu) 
\fancyhead[LO]{\sc \nouppercase{\rightmark}}
% lewa pagina: zawartość \leftmark do prawego, wewnętrznego (marginesu) 
\fancyhead[RE]{\sc \leftmark}

\renewcommand{\chaptermark}[1]{\markboth{\thechapter.\ #1}{}}

% HEAD RULE - IT'S A LINE WHICH SEPARATES HEADER AND FOOTER FROM CONTENT
\renewcommand{\headrulewidth}{0 pt} % 0 MEANS NO RULE, 0.5 MEANS FINE RULE, THE BIGGER VALUE THE THICKER RULE


\fancypagestyle{plain}{
  \fancyhf{}
  \fancyfoot[LE,RO]{\thepage}
  
  \renewcommand{\headrulewidth}{0pt}
  \renewcommand{\footrulewidth}{0.0pt}
}



% --------------------------- CHAPTER HEADERS ---------------------

\usepackage{titlesec}
\titleformat{\chapter}
  {\normalfont\Large \bfseries}
  {\thechapter.}{1ex}{\Large}

\titleformat{\section}
  {\normalfont\large\bfseries}
  {\thesection.}{1ex}{}
\titlespacing{\section}{0pt}{30pt}{20pt} 

    
\titleformat{\subsection}
  {\normalfont \bfseries}
  {\thesubsection.}{1ex}{}


% ----------------------- TABLE OF CONTENTS SETUP ---------------------------

\def\cleardoublepage{\clearpage\if@twoside
\ifodd\c@page\else\hbox{}\thispagestyle{empty}\newpage
\if@twocolumn\hbox{}\newpage\fi\fi\fi}


% THIS MAKES DOTS IN TOC FOR CHAPTERS
\usepackage{etoolbox}
\makeatletter
\patchcmd{\l@chapter}
  {\hfil}
  {\leaders\hbox{\normalfont$\m@th\mkern \@dotsep mu\hbox{.}\mkern \@dotsep mu$}\hfill}
  {}{}
\makeatother

\usepackage{titletoc}
\makeatletter
\titlecontents{chapter}% <section-type>
  [0pt]% <left>
  {}% <above-code>
  {\bfseries \thecontentslabel.\quad}% <numbered-entry-format>
  {\bfseries}% <numberless-entry-format>
  {\bfseries\leaders\hbox{\normalfont$\m@th\mkern \@dotsep mu\hbox{.}\mkern \@dotsep mu$}\hfill\contentspage}% <filler-page-format>

\titlecontents{section}
  [1em]
  {}
  {\thecontentslabel.\quad}
  {}
  {\leaders\hbox{\normalfont$\m@th\mkern \@dotsep mu\hbox{.}\mkern \@dotsep mu$}\hfill\contentspage}

\titlecontents{subsection}
  [2em]
  {}
  {\thecontentslabel.\quad}
  {}
  {\leaders\hbox{\normalfont$\m@th\mkern \@dotsep mu\hbox{.}\mkern \@dotsep mu$}\hfill\contentspage}
\makeatother



% ---------------------- TABLES AD FIGURES NUMBERING ----------------------

\renewcommand*{\thetable}{\arabic{chapter}.\arabic{table}}
\renewcommand*{\thefigure}{\arabic{chapter}.\arabic{figure}}


% ------------- DEFINING ENVIRONMENTS FOR THEOREMS, DEFINITIONS ETC. ---------------

\makeatletter
\newtheoremstyle{definition}
{3ex}%                           % Space above
{3ex}%                           % Space below
{\upshape}%                      % Body font
{}%                              % Indent amount
{\bfseries}%                     % Theorem head font
{.}%                             % Punctuation after theorem head
{.5em}%                          % Space after theorem head, ' ', or \newline
{\thmname{#1}\thmnumber{ #2}\thmnote{ (#3)}}
\makeatother

\theoremstyle{definition}
\newtheorem{theorem}{Theorem}[chapter]
\newtheorem{lemma}[theorem]{Lemma}
\newtheorem{example}[theorem]{Example}
\newtheorem{proposition}[theorem]{Proposition}
\newtheorem{corollary}[theorem]{Corollary}
\newtheorem{definition}[theorem]{Definition}
\newtheorem{remark}[theorem]{Remark}

% --------------------- END OF PREAMBLE PART (MOSTLY) --------------------------





% -------------------------- USER SETTINGS ---------------------------

\newcommand{\tytul}{POLISH TITLE}
\renewcommand{\title}{ENGLISH TITLE}
\newcommand{\type}{Master} % Master OR Engineer
\newcommand{\supervisor}{dr inż. Promotor X} % TITLE AND NAME OF THE SUPERVISOR



\begin{document}
\sloppy
\selectlanguage{english}

\includepdf[pages=-]{title_page} % THIS INPUTS THE TITLE PAGE

\null\thispagestyle{empty}\newpage

% ------------------ PAGE WITH SIGNATURES --------------------------------

%\thispagestyle{empty}\newpage
%\null
%
%\vfill
%
%\begin{center}
%\begin{tabular}[t]{ccc}
%............................................. & \hspace*{100pt} & .............................................\\
%supervisor's signature & \hspace*{100pt} & author's signature
%\end{tabular}
%\end{center}
%


% ---------------------------- ABSTRACTS -----------------------------

{  \fontsize{12}{14} \selectfont
    \begin{abstract}

        \begin{center}
            \title
        \end{center}
        Lorem ipsum dolor sit amet, consetetur sadipscing elitr, sed diam nonumyeirmod tempor invidunt ut labore et dolore magna aliquyam erat, sed diamvoluptua. At vero eos et accusam et justo duo dolores et ea rebum. Stet clita kasd gubergren, no sea takimata sanctus est Lorem ipsum dolor sit amet.

        Lorem ipsum dolor sit amet, consetetur sadipscing elitr, sed diam nonumyeirmod tempor invidunt ut labore et dolore magna aliquyam erat, sed diamvoluptua. At vero eos et accusam et justo duo dolores et ea rebum. Stet clita kasd gubergren, no sea takimata sanctus est Lorem ipsum dolor sit amet.\\

        \noindent \textbf{Keywords:} keyword1, keyword2, ...
    \end{abstract}
}

\null\thispagestyle{empty}\newpage


{\selectlanguage{polish} \fontsize{12}{14}\selectfont
    \begin{abstract}

        \begin{center}
            \tytul
        \end{center}

        Lorem ipsum dolor sit amet, consetetur sadipscing elitr, sed diam nonumyeirmod tempor invidunt ut labore et dolore magna aliquyam erat, sed diamvoluptua. At vero eos et accusam et justo duo dolores et ea rebum. Stet clita kasd gubergren, no sea takimata sanctus est Lorem ipsum dolor sit amet.

        Lorem ipsum dolor sit amet, consetetur sadipscing elitr, sed diam nonumyeirmod tempor invidunt ut labore et dolore magna aliquyam erat, sed diamvoluptua. At vero eos et accusam et justo duo dolores et ea rebum. Stet clita kasd gubergren, no sea takimata sanctus est Lorem ipsum dolor sit amet.\\

        \noindent \textbf{Słowa kluczowe:} slowo1, slowo2, ...
    \end{abstract}
}


%% --------------------------- DECLARATIONS ------------------------------------
%
%%
%%	IT IS NECESSARY OT ATTACH FILLED-OUT AUTORSHIP DEECLRATION. SCAN (IN PDF FORMAT) NEEDS TO BE PLACED IN scans FOLDER AND IT SHOULD BE CALLED, FOR EXAMPLE, DECLARATION_OF_AUTORSHIP.PDF. IF THE FILENAME OR FILEPATH IS DIFFERENT, THE FILEPATH IN THE NEXT COMMAND HAS TO BE ADJUSTED ACCORDINGLY.
%%
%%	command attacging the declarations of autorship
%%
%\includepdf[pages=-]{scans/declaration-of-autorship}
%\null\thispagestyle{empty}\newpage
%
%% optional declaration
%%
%%	command attaching the declaataration on granting a license
%%
%\includepdf[pages=-]{scans/declaration-on-granting-a-license}
%%
%%	.tex corresponding to the above PDF files are present in the 3. declarations folder 
%
\null\thispagestyle{empty}\newpage
% ------------------- TABLE OF CONTENTS ---------------------
% \selectlanguage{english} - for English
\pagenumbering{gobble}
\tableofcontents
\thispagestyle{empty}
\newpage % IF YOU HAVE EVEN QUANTITY OD PAGES OF TOC, THEN REMOVE IT OR ADD \null\newpage FOR DOUBLE BLANK PAGE BEFORE INTRODUCTION


% -------------------- THE BODY OF THE THESIS --------------------------------

\null\thispagestyle{empty}\newpage
\pagestyle{fancy}
\pagenumbering{arabic}
\setcounter{page}{11}


\chapter{Introduction}
\markboth{}{Introduction}
\addcontentsline{toc}{chapter}{Introduction}

What is the thesis about? What is the content of it? What is the Author's contribution to it?
\par
WARNING!  In a diploma thesis which is a team project: Description of the work division in the team, including the scope of each co-author’s contribution to the practical part (Team Programming Project) and the descriptive part of the diploma thesis.
\par

Lorem ipsum dolor sit amet, consetetur sadipscing elitr, sed diam nonumyeirmod tempor invidunt ut labore et dolore magna aliquyam erat, sed diamvoluptua. At vero eos et accusam et justo duo dolores et ea rebum. Stet clita kasd gubergren, no sea takimata sanctus est Lorem ipsum dolor sit amet. Lorem ipsum dolor sit amet, consetetur sadipscing elitr, sed diam nonumyeirmod tempor invidunt ut labore et dolore magna aliquyam erat, sed diamvoluptua. At vero eos et accusam et justo duo dolores et ea rebum. Stet clita kasd gubergren, no sea takimata sanctus est Lorem ipsum dolor sit amet.

\chapter{Example chapter}

This \TeX~file is to be compiled with pdfLaTeX (it's just quick build in TeXMaker).


\section{Example section}

\begin{definition}[Definition]
    %
    A \emph{definition} is a statement of the meaning of a term (a word, phrase, or other set of symbols).
    %
\end{definition}

\subsection{Example subsection}

It's the deepest deph of sectioning allowed by rector.

\begin{definition}[Equation]
    %
    In mathematics, an \emph{equation} is a statement of an equality containing one or more variables.
    %
\end{definition}
%
\begin{example}
    This is an example of an equation:
    \begin{equation}
        %
        2+2=4.
        %
    \end{equation}
    %
    Equation without a number:
    \begin{equation*}
        %
        2+2=4,
        %
    \end{equation*}
    or:
    \[
        2+2=4.
    \]

    It is worthwhile to peruse other mathematical environments like  \emph{multline}, \emph{align} and their versions with a star (, i.e. without numeration). The description of their use can be found at  \url{https://texdoc.org/serve/amsldoc.pdf/0} starting from the end of the third page.

    Equation \eqref{equation} is false. References (and some other things) work properly after compliling \TeX ~ file twice.

    \begin{equation}\label{equation}
        %
        \int \limits_{0}^{1} x \; dx = \frac{3}{2}.
        %
    \end{equation}

\end{example}

Theorem \ref{Pitagoras} is a very interensting result.

\begin{theorem}[Pythagoras' Theorem]\label{Pitagoras}
    %
    Let $c$ represent the length of the hypotenuse and $a$ and $b$ the lengths of the triangle's other two sides. Then:
    %
    \begin{equation*}
        %
        a^2 + b^2 = c^2.
        %
    \end{equation*}
    %
\end{theorem}

\begin{proof}
    %
    The proof has been presented in \cite{Ktos} and \cite{Innyktos}. We can write then \cite{Ktos, Innyktos}.
    %
\end{proof}

\begin{corollary}
    %
    The use of the term \emph{corollary}, rather than \emph{proposition} or \emph{theorem}, is intrinsically subjective.
    %
\end{corollary}




\begin{remark}
    %
    You can find a rather comprehensive list of available symbols at \url{https://www3.nd.edu/~nmark/UsefulFacts/LaTeX_symbols.pdf}.

    If you want to find a symbol by its shape, you can use the following site: \url{https://detexify.kirelabs.org/classify.html}.
    %
\end{remark}

\begin{lemma}[Someone's Lemma]
    %
    Ten lemat jest nie na temat.
    %
\end{lemma}
\begin{proof} Dowód przez indukcję.
\end{proof}


Lorem ipsum dolor sit amet, consetetur sadipscing elitr, sed diam nonumyeirmod tempor invidunt ut labore et dolore magna aliquyam erat, sed diamvoluptua. At vero eos et accusam et justo duo dolores et ea rebum. Stet clita kasd gubergren, no sea takimata sanctus est Lorem ipsum dolor sit amet.Lorem ipsum dolor sit amet, consetetur sadipscing elitr, sed diam nonumyeirmod tempor invidunt ut labore et dolore magna aliquyam erat, sed diamvoluptua. At vero eos et accusam et justo duo dolores et ea rebum. Stet clita kasd gubergren, no sea takimata sanctus est Lorem ipsum dolor sit amet.



\section{Floats -- tables and figures}

\begin{table}% label has to be after caption, otherwise numeration is wrong
    \caption[Short caption]{Additional options}
    \label{opcje}
    \centering
    \begin{tabular}{|c|p{0.8\textwidth}|}
        \hline
        symbol     & effect                                                                                                                                                                                 \\ \hline
        \texttt{h} & Place the float here, i.e., approximately at the same point it occurs in the source text (however, not exactly at the spot)                                                            \\
        \texttt{t} & Position at the top of the page                                                                                                                                                        \\
        \texttt{b} & Position at the bottom of the page                                                                                                                                                     \\
        \texttt{p} & Put on a special page for floats only                                                                                                                                                  \\
        \texttt{!} & Override internal parameters LaTeX uses for determining "good" float positions                                                                                                         \\
        \texttt{H} & Places the float at precisely the location in the \LaTeX ~ code. Requires the float package,[1] i.e., \texttt{\textbackslash usepackage\{float\}}. This is somewhat equivalent to !ht. \\ \hline
    \end{tabular}
\end{table}

Place labels after captions or you get the wrong labelling.

In Table \ref{opcje} there are additional options for \texttt{table} and \texttt{figure} environments.

\begin{figure}[h!]

    \begin{center}
        \setlength{\unitlength}{1mm}

        \begin{picture}(40, 30)
            \put(20,1){\line(0,1){20}} % line

            % bottom
            \put(20,1){\circle*{2}}
            \put(25,1){0}

            % top
            \put(20,21){\circle*{2}}
            \put(25,21){1}
        \end{picture}

    \end{center}
    \caption{Example figure -- it has been drawn by \LaTeX ~default tools}
\end{figure}


Lorem ipsum dolor sit amet, consetetur sadipscing elit, sed diam nonumyeirmod tempor invidunt ut labore et dolore magna aliquyam erat, sed diamvoluptua. At vero eos et accusam et justo duo dolores et ea rebum. Stet clita kasd gubergren, no sea takimata sanctus est Lorem ipsum dolor sit amet.Lorem ipsum dolor sit amet, consetetur sadipscing elitr, sed diam nonumyeirmod tempor invidunt ut labore et dolore magna aliquyam erat, sed diamvoluptua. At vero eos et accusam et justo duo dolores et ea rebum. Stet clita kasd gubergren, no sea takimata sanctus est Lorem ipsum dolor sit amet.




\chapter{The next chapter}

Lorem ipsum dolor sit amet, consetetur sadipscing elit, sed diam nonumyeirmod tempor invidunt ut labore et dolore magna aliquyam erat, sed diamvoluptua. At vero eos et accusam et justo duo dolores et ea rebum. Stet clita kasd gubergren, no sea takimata sanctus est Lorem ipsum dolor sit amet.Lorem ipsum dolor sit amet, consetetur sadipscing elitr, sed diam nonumyeirmod tempor invidunt ut labore et dolore magna aliquyam erat, sed diamvoluptua. At vero eos et accusam et justo duo dolores et ea rebum. Stet clita kasd gubergren, no sea takimata sanctus est Lorem ipsum dolor sit amet.


\section{Matrices}

Simple matrix:
\begin{equation*}
    \begin{matrix}
        a & b & c & d \\
        d & e & f & g \\
        1 & 1 & 1 & 1
    \end{matrix}
\end{equation*}
%
Matrix with parentheses:
%
\begin{equation*}
    A =
    \begin{pmatrix}
        a & b & c & d \\
        d & e & f & g \\
        1 & 1 & 1 & 1
    \end{pmatrix}
\end{equation*}
%
Matrix with brackets:
%
\begin{equation*}
    \begin{bmatrix}
        a & b & c & d \\
        d & e & f & g \\
        1 & 1 & 1 & 1
    \end{bmatrix}
\end{equation*}
%
You can also use more general environment:
%
\begin{equation*}
    \renewcommand{\arraystretch}{0.8}
    \begin{array}{ccc}
        1 & 0 & 0 \\
        0 & 1 & 0 \\
        0 & 0 & 1 \\
    \end{array}
\end{equation*}
%
Matrix with braces:
%
\begin{equation*}
    \left\{
    \renewcommand{\arraystretch}{0.8}
    \begin{array}{ccc}
        1 & 0 & 0 \\
        0 & 1 & 0 \\
        0 & 0 & 1 \\
    \end{array}\right\}
\end{equation*}

\begin{definition}
    %
    Let $A\neq \emptyset$, $n \in \mathbb{N}$. Every function $f\colon A^n \to A$ is called an \emph{$n$-ary operation} or \emph{działaniem} określonym na $A$.
    $0$-ary operations are constant functions.
    %
\end{definition}


\begin{definition}[Algebra]
    %
    The ordered pair $(A,F)$, where $A\neq \emptyset$ is a set and $F$ is a family of operations defined on $A$, shall be called an \emph{algebra} (or \emph{$F$-algebra}). The set $A$ is called \emph{the set of elements}, \emph{support} or \emph{universe} of an algebra $(A,F)$ and $F$ is called \emph{the set of elementary operations}.
    %
\end{definition}

\begin{proposition}
    %
    I state that, having passed to the limit, the only thing left me me is to camp at said limit or return, or, maybe, search for a pass or an exit to other areas.
    %
\end{proposition}





% ------------------------------- BIBLIOGRAPHY ---------------------------
% LEXICOGRAPHICAL ORDER BY AUTHORS' LAST NAMES
% FOR AMBITIOUS ONES - USE BIBTEX


\begin{thebibliography}{20} % IF YOU HAVE MORE REFERENCES, WRITE THE BIGGER NUMBER

    \bibitem[1]{Ktos} A. Author, \emph{Title of a book}, Publisher, year, page--page.
    \bibitem[2]{Innyktos} J. Bobkowski, S. Dobkowski, Title of an article, \emph{Magazine X, No. 7}, year, PAGE--PAGE.
    \bibitem[3]{B} C. Brink, Power structures, \emph{Algebra Universalis 30(2)}, 1993, 177--216.
    \bibitem[4]{H} F. Burris, H. P. Sankappanavar, \emph{A Course of Universal Algebra}, Springer-Verlag, New York, 1981.
\end{thebibliography}
\pagenumbering{gobble}
\thispagestyle{empty}



% ----------------------- LIST OF SYMBOLS AND ABBREVIATIONS ------------------
\chapter*{List of symbols and abbreviations}

\begin{tabular}{cl}
    nzw.           & nadzwyczajny  \\
    *              & star operator \\
    $\widetilde{}$ & tilde
\end{tabular}
\\
If you don't need it, delete it.
\thispagestyle{empty}


% ----------------------------  LIST OF FIGURES --------------------------------
\listoffigures
\thispagestyle{empty}
If you don't need it, delete it.


% -----------------------------  LIST OF TABLES --------------------------------
\renewcommand{\listtablename}{Spis tabel}
\listoftables
\thispagestyle{empty}
If you don't need it, delete it.

% -----------------------------  LIST OF APPENDICES ---------------------------
\chapter*{List of appendices}
\begin{enumerate}
    \item Appendix 1
    \item Appendix 2
    \item In case of no appendices, delete this part.
\end{enumerate}
\thispagestyle{empty}


\end{document}
